% Authored. Tables and numbers come from generated/.

\section{A European Benchmark: Is This Composition Unusual?}
\label{sec:benchmark}

Everything to this point describes one country. It establishes how Portugal's
general-government balance is composed and how persistent that composition is; it
cannot establish whether the composition is unusual, because a single-country study
has no comparison set. That distinction matters for how the results should be read:
a subsector pattern that is common across Europe invites a different explanation from
one that is peculiar to Portugal.

ESA 2010 requires the same subsector breakdown from every reporter, so the comparison
is available \citep{eurostat_gov_main,ecb_gfs}. We take net lending and net borrowing
by subsector for 1995--\PanelEnd\ from Eurostat and repeat the domestic definitions on
it.

\subsection{Construction}

Three choices matter, and each is a place where a careless benchmark would mislead.

\textbf{The non-Social-Security aggregate must include state government.} Portugal has
no S.1312 tier, so its aggregate is central plus local. Germany, Spain, Austria,
Belgium and Switzerland do have one, and omitting it would leave their identity open
and understate their non-Social-Security deficits. We therefore define
$\Bnonssf_{c,t} = B^{S.1311}_{c,t} + B^{S.1312}_{c,t} + B^{S.1313}_{c,t}$, treating a
missing state tier as contributing nothing because the tier does not exist rather than
because a value is unknown. With that definition the identity closes to published
rounding for every reporter.

\textbf{Ratios are computed in national currency.} Eurostat publishes shares of GDP to
one decimal, which is an unusable denominator for a ratio: a non-Social-Security
balance printed as $-0.2$ could lie anywhere in a band wide enough to move the offset
ratio by a quarter of its value. The levels carry far more precision, and the ratio is
currency-invariant. We additionally require the denominator to be at least $0.5$\% of
GDP in absolute size before reporting a ratio, so that a near-zero denominator cannot
manufacture a large one.

\textbf{Coverage is made comparable.} A country-year enters only if the four required
sectors --- the aggregate, central government, local government and social security funds
--- are all reported. The state tier is required conditionally rather than absolutely: a
country that reports S.1312 anywhere in its record must report it in that year too, while
a country that never reports it, Portugal among them, is not excluded for lacking a tier
it does not have. Requiring the tier unconditionally would drop every unitary reporter;
not requiring it at all would admit federal country-years whose non-Social-Security
aggregate is incomplete, and it is exactly those years in which the identity would fail
to close. A reporter then enters the summary only with at least fifteen complete years,
because a frequency computed over a handful of years is not comparable with one computed
over thirty. This leaves \BenchmarkCountries\ reporters over
\BenchmarkStart--\BenchmarkEnd, with Portugal contributing \BenchmarkYears\ years.

Two remarks on what the comparison is worth. Eurostat's Portuguese rows agree with our
domestic panel to rounding, which is an independent check on the extraction of
Section~\ref{sec:data} as well as the basis for the comparison. And this is a
distribution, not a test: locating Portugal within a spread of accounting compositions
says how common that composition is, not why any country's takes the form it does.

\subsection{Where Portugal sits}

Table~\ref{tab:benchmark} reports the composition for every reporter and
Table~\ref{tab:benchmarkposition} locates Portugal in each distribution.

\begin{table}[htbp]
\centering
\caption{Subsector composition across reporters, ordered by the frequency of a Social Security surplus. Sign frequencies are shares of the years each reporter covers completely.}
\label{tab:benchmark}
\small
\fitwidth{%
\begin{tabular}{@{}lrrrrrrr@{}}
\toprule
Reporter & N & Central $<0$ & SSF $>0$ & Mean SSF (\% GDP) & Surplus years & of which non-SSF $<0$ & Median offset \\
\midrule
FI & 31 & 0.774 & 1.000 & 2.45 & 11 & 6 & 0.540 \\
LU & 31 & 0.645 & 1.000 & 1.65 & 24 & 10 & 0.892 \\
DK & 31 & 0.355 & 0.935 & 0.05 & 21 & 0 & 0.008 \\
CY & 31 & 0.806 & 0.935 & 2.23 & 9 & 4 & 0.446 \\
PT & 31 & 1.000 & 0.935 & 0.71 & 4 & 4 & 0.141 \\
IT & 31 & 1.000 & 0.839 & 0.18 & 0 & 0 & 0.064 \\
EL & 31 & 0.903 & 0.839 & 0.97 & 6 & 3 & 0.206 \\
EE & 31 & 0.613 & 0.839 & 0.25 & 12 & 2 & 0.306 \\
IS & 28 & 0.536 & 0.821 & 0.16 & 10 & 0 & 0.043 \\
SE & 31 & 0.581 & 0.806 & 0.19 & 14 & 2 & 0.225 \\
SK & 31 & 1.000 & 0.742 & 0.42 & 0 & 0 & 0.061 \\
BG & 31 & 0.613 & 0.710 & 0.17 & 13 & 2 & 0.095 \\
RO & 31 & 1.000 & 0.710 & 0.13 & 0 & 0 & 0.084 \\
CH & 30 & 0.533 & 0.700 & 0.16 & 17 & 5 & 0.058 \\
LV & 31 & 0.968 & 0.677 & 0.35 & 2 & 0 & 0.316 \\
LT & 31 & 1.000 & 0.613 & 0.25 & 4 & 3 & 0.419 \\
DE & 31 & 0.774 & 0.613 & 0.04 & 8 & 2 & 0.085 \\
NL & 31 & 0.806 & 0.581 & 0.17 & 8 & 4 & 0.310 \\
IE & 31 & 0.516 & 0.581 & 0.09 & 16 & 3 & 0.167 \\
BE & 31 & 1.000 & 0.581 & 0.06 & 3 & 2 & 0.072 \\
AT & 31 & 0.968 & 0.548 & 0.02 & 2 & 0 & 0.056 \\
HU & 31 & 1.000 & 0.548 & 0.02 & 0 & 0 & 0.045 \\
SI & 31 & 0.935 & 0.516 & -0.05 & 3 & 1 & 0.043 \\
PL & 31 & 1.000 & 0.484 & 0.24 & 0 & 0 & 0.263 \\
HR & 31 & 0.871 & 0.452 & -0.12 & 4 & 1 & 0.124 \\
FR & 31 & 1.000 & 0.419 & -0.20 & 0 & 0 & 0.072 \\
CZ & 31 & 0.935 & 0.387 & -0.04 & 4 & 0 & 0.051 \\
ES & 31 & 0.903 & 0.355 & -0.27 & 3 & 0 & 0.464 \\
\bottomrule
\end{tabular}%
}
\par\smallskip
\begin{minipage}{\linewidth}\footnotesize\itshape The non-Social-Security aggregate is central plus state plus local government, so federal reporters are treated consistently with unitary ones. Reporters with fewer than fifteen complete years are excluded, because a frequency over a handful of years is not comparable with one over thirty.\end{minipage}
\end{table}

\begin{table}[htbp]
\centering
\caption{Portugal's position in each cross-country distribution. The percentile is the share of reporters below Portugal's value.}
\label{tab:benchmarkposition}
\small
\fitwidth{%
\begin{tabular}{@{}lrrrrr@{}}
\toprule
Metric & Portugal & Cross-country median & Minimum & Maximum & Percentile \\
\midrule
Share of years with a Central Government deficit & 1.000 & 0.903 & 0.355 & 1.000 & 68 \\
Share of years with a Social Security surplus & 0.935 & 0.689 & 0.355 & 1.000 & 82 \\
Share of years with an aggregate surplus & 0.129 & 0.129 & 0.000 & 0.774 & 39 \\
Mean Social Security balance (\% GDP) & 0.706 & 0.168 & -0.271 & 2.448 & 82 \\
Median offset ratio, where defined & 0.141 & 0.110 & 0.008 & 0.892 & 54 \\
\bottomrule
\end{tabular}%
}
\end{table}


The answer is not uniform across the paper's findings, and that is the substantive
result of this section.

\paragraph{The Central Government deficit is not unusual.} Portugal records a Central
Government deficit in \PtCentralNegativeShare\% of its reported years, which places it
at the \PtCentralPercentile th percentile --- but \CentralAlwaysNegativeCountries\ of
the \BenchmarkCountries\ reporters also record one in every single year. A persistently
deficit-running central tier is a common European pattern, not a Portuguese
peculiarity. Read on its own, the domestic finding that Central Government is negative
in all \CentralNegativeYears\ years would invite the opposite impression.

\paragraph{The Social Security surplus is unusual, in both frequency and size.}
Portugal records a Social Security surplus in \PtSsfPositiveShare\% of years against a
cross-country median of \PtSsfShareMedian\%, at the \PtSsfSharePercentile th
percentile. Its mean Social Security balance is \PtSsfMeanPct\% of GDP against a median
of \PtSsfMeanMedian\%, again at the \PtSsfMeanPercentile th percentile. On both
measures Portugal sits in the upper tail.

\paragraph{The typical offset magnitude is ordinary.} Portugal's median offset ratio,
over the years in which it is defined, is \PtOffsetMedian\ against a cross-country
median of \OffsetMedianAcrossCountries, at the \PtOffsetPercentile th percentile. In a
typical year the Portuguese Social Security surplus offsets no more of the
non-Social-Security deficit than elsewhere. What distinguishes Portugal is the recent
years, in which the ratio exceeds one, not the central tendency.

That conclusion rests on the denominator floor, which is a choice rather than a property
of the data, so it is worth checking whether the conclusion survives other choices.
Sweeping the floor from \OffsetFloorLow\ to \OffsetFloorHigh\ percent of GDP moves
Portugal's percentile only between \OffsetPercentileLow\ and \OffsetPercentileHigh: the
position stays near the middle of the distribution throughout. Had the percentile moved
sharply with the floor, the finding would have been an artefact of where we set it.

\begin{figure}[htbp]
\centering
\includegraphics[width=\linewidth]{17_european_benchmark.png}
\caption{Subsector sign frequencies across reporters, Portugal highlighted, ordered by
the frequency of a Social Security surplus. Left: share of reported years with a Social
Security surplus. Right: share with a Central Government deficit. Portugal is in the
upper tail of the first and shares the maximum of the second with several reporters.}
\label{fig:benchmark}
\end{figure}

\begin{figure}[htbp]
\centering
\includegraphics[width=\linewidth]{18_european_offset_distribution.png}
\caption{Offset ratio where defined: Portugal against all other reporters pooled, as a
share of each group's defined country-years. Values above three are clipped to the last
bin. The Portuguese distribution is not shifted relative to the others; its recent years
sit in the right tail of both.}
\label{fig:offsetdistribution}
\end{figure}

\subsection{The composition of a surplus year}

The sharpest comparison is not a sign frequency but the structure of the surplus years
themselves, because that is the paper's headline composition restated as a
cross-country question: when a country records an aggregate surplus, is its
non-Social-Security balance in deficit?

Across the \CountriesWithSurplus\ reporters with at least one aggregate surplus, that
composition holds in \SurplusYearsOffsetting\ of \SurplusYearsTotal\ surplus
country-years, or \SurplusOffsettingShare\%. It is therefore a minority pattern: most
European surpluses are recorded with the non-Social-Security tiers also in surplus, or
close to balance.

Portugal is the exception in the strict sense. Of the \CountriesWithSurplus\ reporters,
\CountriesAllOffsetting\ record a negative non-Social-Security balance in \emph{every}
one of their aggregate-surplus years, and Portugal is among them, at
\PtSurplusOffsetting\ of \PtSurplusYears.

The pooled figure weights by country-year, so a reporter with two dozen surplus years
counts for more than one with three. Weighting by country instead asks a different and
equally reasonable question --- of a country's own surplus years, what share pair with a
negative non-Social-Security balance? --- and the two agree: the country-weighted median
is \SurplusCountryMedian\%, with quartiles at \SurplusCountryLowerQuartile\% and
\SurplusCountryUpperQuartile\%, against the pooled \SurplusOffsettingShare\%. Portugal
sits at the \SurplusCountryPercentile th percentile of that distribution. Because the two
weightings can disagree and neither is obviously correct, both are reported.

That result needs its sample stated plainly alongside it. Portugal has only
\PtSurplusYears\ surplus years, so ``all of them'' is four observations, and several
reporters with more surplus years have high but not unanimous shares. The defensible
statement is therefore the conjunction: the composition is a minority pattern in Europe,
Portugal exhibits it in every surplus year it has recorded, and the number of such years
is small. We would not read a stronger claim than that from
\PtSurplusYears\ observations.

\subsection{What the benchmark changes}

It changes the reading of two results and confirms a third.

The Central Government finding is weaker than it looks in isolation: a permanently
deficit-running central tier is normal in Europe. The Social Security finding is
stronger: both the frequency and the size of the Portuguese surplus sit in the upper
tail. And the composition of the surplus years is genuinely distinctive, subject to the
sample caveat above.

None of this identifies a cause. Countries differ in how they assign contributory
schemes between tiers, in whether they operate a state-government level, in the maturity
of their pension systems and in how they route transfers between subsectors --- and the
benchmark holds none of that constant. It establishes that the Portuguese composition is
unusual in a specific, stated respect, which is what the rest of the paper could not do.
